\begingroup
\begin{figure*}[ht]
\centering
\hspace{-.1in}
	\begin{subfigure}[b]{0.32\textwidth}
		\includegraphics[width=\textwidth]{A_PBQR_CURVATUREgms.png}
		\subcaption[]{Compute}
		\label{fig:A_PBQR_CURVATUREgms}
	\end{subfigure}
\hspace{-.1in}
	\begin{subfigure}[b]{0.32\textwidth}
		\includegraphics[width=\textwidth]{A_PBQR_CURVATUREcms.png}
		\subcaption[]{Graphics}
		\label{fig:A_PBQR_CURVATUREcms}
	\end{subfigure}
\hspace{-.1in}
	\begin{subfigure}[b]{0.32\textwidth}
		\includegraphics[width=\textwidth]{A_PBQR_CURVATUREtot.png}
		\subcaption[]{Total}
		\label{fig:A_PBQR_CURVATUREtot}
	\end{subfigure}
\hspace{-.1in}
\caption[TITLE:PQB Curvature Bound Analysis]{Plots of quadratic fit lines for the saturated range of \numprint{\APBQRCURVATUREStartRowParticlesVal{}} to \numprint{\APBQRCURVATUREgmsMaxParticles{}} particles. The curvature for the graphics pipeline is \APBQRCURVATUREgmsa{}$N^2$. The curvature for the compute pipeline is \APBQRCURVATUREcmsa{}$N^2$ indicating a negative curvature. The total  curvature is \APBQRCURVATUREtota{}$N^2$. The total quadratic contribution is \APBQRCURVATUREtotQuadraticContributionpct{} percent.}
\Description[TITLE:PQB Curvature Bound Analysis]{Plots of quadratic fit lines for the saturated range of \numprint{\APBQRCURVATUREStartRowParticlesVal{}} to \numprint{\APBQRCURVATUREgmsMaxParticles{}} particles. The curvature for the graphics pipeline is \APBQRCURVATUREgmsa{}$N^2$. The curvature for the compute pipeline is \APBQRCURVATUREcmsa{}$N^2$ indicating a negative curvature. The total  curvature is \APBQRCURVATUREtota{}$N^2$. The total quadratic contribution is \APBQRCURVATUREtotQuadraticContributionpct{} percent.}
\label{fig:A_PBQR_CURVATURE}
\end{figure*}
\endgroup